\chapter{Die sechs neuen Quellen: geprüfte Korrekturen}
\label{ch:audit13}
\section{Was gegenüber v1.1 neu ist}
Die sechs zusätzlichen Unterlagen heißen im neuen Quellenregister N1 bis N6; die ursprünglichen S01 bis S11 bleiben unverändert. Ihr Audit trennt exakte endliche Identitäten, numerische Spektren und zusätzliche Modellannahmen. 100 eigene Bedingungen bestanden: 78 exakte beziehungsweise Integritätsprüfungen und 22 numerische. Die anschließende gemeinsame Ausführung hat einen eigenen 231-Bedingungen-Vertrag. Diese Zählungen messen keine unabhängigen Entdeckungen.

\section{Die zentrale Reparatur: Wege nicht zu früh unterscheiden}
Setze $S\ket{ab}=\ket{ba}$ und $D=\sum_a\ket{aa}\bra{aa}$. Der Wedge-Kanal besitzt
\begin{equation}
K\ket{ab}=\ket{a\wedge b},\quad K\ket{ba}=-\ket{a\wedge b},\quad
K^\dagger K=I-S,\quad KK^\dagger=2I_6.
\end{equation}
Für $a=b$ verschwindet $K$. Beide geordneten Wege enden im gleichen virtuellen Zustand und interferieren. Orthogonale Records für $ab$ und $ba$ ersetzen dies durch $\widetilde K^\dagger\widetilde K=I-D$: Der antisymmetrische Rang-6-Kanal wird zum Rang-12-Kanal. Bei gleichem Energienenner entsteht eine Kollisionsenergie statt des gewünschten $\mathrm{SU}(4)$-Austauschs. Auf dem vollständigen Vierergraphen bleiben 24 Farbanordnungen im Grundraum; $\Omega$ ist nicht eindeutig.

Die Korrektur lautet $h_{e,ab}=h_{e,ba}$ innerhalb derselben Kante. Verschiedene Kanten dürfen durch Loch- und Zuschauerzustände unterscheidbar bleiben. Für reellen Überlapp $0\le\eta\le1$ lautet die interpolierende Gram-Matrix
\begin{equation}
G_\eta=I-\eta S-(1-\eta)D.
\end{equation}
Nach gemeinsamer Energieverschiebung erhält man
\begin{equation}
\overline H_\eta=2\eta H_{\mathrm{tet}}/J+(1-\eta)D_{\mathrm{coll}},
\qquad D_{\mathrm{coll}}=\sum_{i<j}D_{ij}.
\end{equation}
Für $\eta>0$ ist $\Omega$ eindeutig und die dimensionslose Lücke genau $4\eta$. Die untere Schranke folgt aus der Tetramerlücke und Positivität; sie wird im verschiedenen-Farben-Sektor erreicht. Führend entspricht dies $4\eta t^2/\Delta$. Bei $\eta=0$ springt die Kerndimension auf 24; die positive Lücke über dem gesamten Kern ist eins. Ein Elternoperator aus dem Projektor auf $\ker K_\eta$ wäre eine andere, diskontinuierliche Konstruktion.

\begin{bild}{Bildlich}
Zwei Wege dürfen ein gemeinsames Ziel erreichen, ohne ein unterscheidbares Wegschild zu hinterlassen. Dieses Schild würde ihre Interferenz zerstören. Aufzeichnen lässt sich stattdessen, ob ein Vermittlungsereignis stattfand. Diese gröbere Information verwendet die neue gemeinsame Ausführung.
\end{bild}

\section{Der erste Schatten wählt die Fortsetzung noch nicht}
Mit $P_\pm=(I\pm S)/2$ und $W=K/\sqrt2$ ist
\begin{equation}
U_\theta=\begin{pmatrix}e^{i\theta}P_+&W^\dagger\\W&0\end{pmatrix}
\end{equation}
auf $(\C^4\otimes\C^4)\oplus\C^6$ unitär und kollektiv $\mathrm{SU}(4)$-kovariant. Alle $\theta$ liefern für Materieeingänge denselben ersten blockdephasierten Schatten. Dennoch wirkt $U_0^2$ dort als $I$, während $U_{\pi/2}^2$ als $-S$ wirkt: $\ket{01}$ kehrt als $\ket{01}$ oder als $-\ket{10}$ zurück. Das Nullereignis unverändert zu lassen wählt $\theta=0$ in dieser Klasse, ist aber ein zusätzliches Prinzip.

\section{Der allkopplige Zweizellenbeweis}
Für genau zwei vollständige $K_4$-Zellen und eine positive Brücke gelten die Formeln aus Kapitel~\ref{ch:two} bei allen endlichen $\lambda\ge0$. Die lokalen Zellenergien sind $0,2J,3J,4J,6J$. Unter $4J$ können global nur das Produktsingulett, adjungierte Sektoren bei $2J$ und die 20 bei $3J$ auftreten. Der bekannte Singulettblock erreicht $E_0$; sein orthogonales Singulettkomplement beginnt bei mindestens $4J$.

Im adjungierten Sektor sei $P$ die $H_0$-Energie-$2J$-Projektion. Dann gilt $H_0\ge4JI-2JP$. Die exakt berechnete Kompression von $V=P^+_{ab}$ besitzt pro adjungierter Komponente
\begin{equation}
\spec(PVP|_{\operatorname{ran}P})
=\{(1/2)^{[1]},(5/8)^{[4]},(3/4)^{[1]}\}.
\end{equation}
Zum Eigenwert $p$ liefert der Zwei-Projektionen-Vergleich den unteren Wert
\begin{equation}
3J+\lambda/2-\tfrac12\sqrt{4J^2+4J\lambda(1-2p)+\lambda^2}.
\end{equation}
Wegen $p\ge1/2$ liegt dieser mindestens bei $E_1=3J+\lambda/2-\sqrt{4J^2+\lambda^2}/2$. Die Schranke wird im tatsächlich invarianten Block
\begin{equation}
\begin{pmatrix}2J+\lambda/2&\lambda/2\\\lambda/2&4J+\lambda/2\end{pmatrix}
\end{equation}
erreicht. Alle anderen Darstellungen beginnen bei mindestens $3J$, während $E_1<3J$. Mit $R$ und $Q$ aus Kapitel~\ref{ch:two} gilt
\begin{equation}
R^2-(Q-J)^2=11J^2+2J(Q-\lambda)>0,\quad Q-J>0.
\end{equation}
Also $R>Q-J$ und $E_1-E_0>J/2$. Der Grundzustand bleibt eindeutig. Für $\lambda\to\infty$ nähern sich $E_0$ und $E_1$ den Werten $5J/2$ und $3J$.

Ganzzahlige Tensoridentitäten wurden auf 65536 Komponenten geprüft; zusätzlich alle 15 zulässigen Youngsektoren bei zehn Kopplungswerten numerisch diagonalisiert. Der allgemeine Satz folgt aus dem Vergleichsbeweis, nicht aus diesen Stichproben. Eine größenuniforme Vielzellenlücke folgt nicht.

\section{Vierte Ordnung und unterschiedliche Subtraktionskonventionen}
Im ausgewiesenen harmonischen Modell eines geteilten Vermittlers lautet der symmetrische aktive Block
\begin{equation}
M_s=\begin{pmatrix}0&\sqrt2g&0\\\sqrt2g&\Delta&2g\\0&2g&2\Delta\end{pmatrix}.
\end{equation}
Sein niedriger Zweig beginnt mit $-2g^2/\Delta+O(g^6/\Delta^5)$; der antisymmetrische mit $-2g^2/\Delta+4g^4/\Delta^3+O(g^6/\Delta^5)$. Deshalb ist auf dem aktiven $6\otimes6$-Raum
\begin{equation}
H^{(4)}_{\mathrm{gesamt}}=4g^4P_-^{(6)}/\Delta^3.
\end{equation}
Nach Abzug beider Einzelbondkorrekturen $2g^4I/\Delta^3$ bleibt
\begin{equation}
H^{(4)}_{\mathrm{verbunden}}=-2g^4S_6/\Delta^3
=-8t^4S_6/\Delta^3,\quad g=\sqrt2t.
\end{equation}
Beide Schreibweisen sind kompatibel. Eine Konstante auf dem aktiven Sektor ist außerhalb davon nicht automatisch global skalar. Der vollständige Clebsch-Operator wurde in dieser früheren Auditstufe noch nicht geprüft. Die anschließenden neuen Eingaben und ihre Gegenprüfung in Kapitel~\ref{ch:newinputs13} erweitern genau diese Stelle.

Im isolierten Kanal ist $j=(\sqrt{\Delta^2+4g^2}-\Delta)/2$ exakt; $2t^2/\Delta$ ist nur führend. Das Verhältnis des verbundenen Koeffizienten zum führenden $J$ ist $4(t/\Delta)^2$: vier Prozent bei $t/\Delta=0{,}1$, ein Prozent bei $0{,}05$. Starker kohärenter Volltransfer ist kein Schwachkopplungsgrenzfall.

\section{Ein exakter Stern-Spektralfilter unter zusätzlichen Ressourcen}
Für $H_\star=J(P^+_{01}+P^+_{02}+P^+_{03})$ gilt
\begin{equation}
\spec(H_\star/J)=
\{0^{[1]},(1/2)^{[30]},1^{[45]},(3/2)^{[40]},2^{[15]},(5/2)^{[90]},3^{[35]}\}.
\end{equation}
Mit $M=2H_\star/J$ und $U=\exp[-i\pi H_\star/(2J)]$ folgen
\begin{equation}
\prod_{j=1}^{6}(jI-M)=720P_\Omega,\qquad
\frac18\sum_{k=0}^{7}U^k=P_\Omega.
\end{equation}
Ein kohärentes Achtzweigregister liefert aus dem früheren Zweisingulett-Eingang $\chi_0$ Erfolg $1/6$. Alle acht Ausgänge haben Gewichte $(1/6,1/6,0,0,1/6,1/6,1/6,1/6)$. Dies ist eine exakte Alternative zur alten 32-Zweigschaltung. Kontrollierte kontinuierliche Sternpulse und Register bleiben zusätzliche Ressourcen. Erzeugt man $\chi_0$ erst aus $\ket{0123}$ durch zwei Paarfilter, kostet das $1/4$; die Gesamtausbeute ist $1/24$.

\section{Einfachheit braucht einen Auslesevertrag}
In der Inzidenzfamilie $K_a$ sind die beiden relevanten Entropien
\begin{align}
H_C(a)&=-a\log a-(1-a)\log((1-a)/6),\\
H_R(a)&=-a\log a-(1-a)\log((1-a)/12).
\end{align}
Strikte Konkavität liefert verschiedene Maxima: $a=1/7$ für Kontexte und $a=1/13$ für Strahlen. Letzterer Prozess hat Rang 60. Erweitert man den Nachfolgerträger, entstehen symmetrieverträgliche Ränge 15 und 16. Weder maximale Entropie noch minimale Rangzahl wählt die Regel ohne erklärte Klasse. Die passive separate Marginalauslesung bleibt bei 45 linearen beziehungsweise 44 normierten affinen Größen; selektive Eingriffe gehören nicht dazu.

\section{Weitere architektonisch wichtige Korrekturen}
\begin{description}
\item[Clebsch.] Zwei Singulettläufe finden mindestens vier orthogonale Moden bei $11{,}561762122803J$ über $11{,}045398337068J$, Residuen unter $10^{-8}$. Das Dublett ist widerlegt; ein exakter oberer Multiplizitätsbeweis fehlt.
\item[$E_8$-Typen.] 960 gleichgradige Wurzelsummen gehen nach $(10,6)$; gemischt ergeben sich 640 und 192 Summen nach Grad null sowie 64 Cartanfälle. Der vorgeschlagene gemischte Kanal nach $(10,6)$ ist typwidrig. Die Korrektur schließt kein natives Half-Charge-Feld.
\item[Ort versus innere Symmetrie.] Im gewählten CAR-Modell liefert Gewichtsplatzmischung $[H_U,T]P=UTP\ne0$. Fester Gewichtsplatz als Ort und unverändert ungebrochene innere $\mathrm{Spin}(10)$-Mischung sind verschiedene Rollen.
\item[Kommutanten.] Die volle Kommutante eines einzelnen Tetramer-Hamiltonoperators hat Dimension $1^2+45^2+40^2+135^2+35^2=23076$. Rekonstruktion verlangt die gesamte lokale Operatorfamilie und deren erlaubte Automorphismen.
\item[Dimension.] Der Clebsch-Graph hat 16 Knoten, 40 Kanten und 25 unabhängige Zyklen. Mit reduzierter Inzidenzmatrix $D_\Gamma$ und Kantenspannungen $B_d$ wurde für vier periodische Überlagerungen in $d=1,2,3,4$ die akustische Form
\[
B_d^T[I-D_\Gamma(D_\Gamma^TD_\Gamma)^{-1}D_\Gamma^T]B_d/16
\]
rational positiv definit geprüft. Dieselbe lokale Inzidenz wählt daher keine Dimension.
\item[Weyl.] Richtig ist $\det(\omega I-\sum_i v_iq_i\sigma_i)=\omega^2-\sum_i v_i^2q_i^2$, ohne die falschen Kreuzterme einer voll summierten Form $v_iv_jq_iq_j$. Allgemein steht dort $\mathbf q^TV^TV\mathbf q$.
\item[Dunkler Zustand.] Für $L=I-P_\Omega$, $H=0$ ist jeder Zustand im orthogonalen Komplement stationär, obwohl der dunkle Raum nur $\C\Omega$ ist. Eindeutiger dunkler Vektor bedeutet nicht globaler Attraktor.
\end{description}
Alpha- und Inflationskorrekturen stehen integriert in Kapitel~\ref{ch:zahlen}. Negative Lepton-, Higgs- und Protonzweige bleiben erhalten. Kein RH-, Faktorisierungs-, P-vs-NP- oder TOE-Abschluss folgt aus dem Audit.
